\xmpheader 3/{Fonts and special characters}
\chardef \\ = `\\ % Let \\ denote a backslash

\font\tencmss = cmss10  % Load some unstandard fonts
\font\tencmssbx = cmssbx10
\font\bigbf = cmbx10 scaled \magstep2  % You can use \magstep1 upto \magstep5.

{\it Here are a few words in an italic font}, {\bf a few words in a boldface font}, {\it and a\/
{\bf mixture} of two, with two\/ {\rm roman words} inserted}.  Where an italic font is followed by
a nonitalic font, we've inserted an ``italic correction'' ({\tt \\/}) to make spacing look right.
Here's a {\sevenrm smaller} word---but the standard \TeX\ fonts won't give you anything smaller
than {\fiverm this}. The standard font like the current one is called cmr10.  {\tencmss But now
we are using cmss10} {\tencmssbx or in its bold version}.  Here is a {\bigbf big A}.  You can magnify 
the whole document by saying {\tt \\magnification \\magstep1}, at the very beginning of the file.

If you need any of the characters:
\medskip
\centerline{\$ \quad \& \quad \# \quad \_ \quad \% \quad 
\char`\^ \quad \char`\~ \quad $\{$ \quad $\}$ \quad $\backslash$}
% The \quad inserts an em space between charaters.
\medskip
\noindent you'll need to write them a special way. Look at the input to see how to do it. 

\TeX\ has the accents and letters that you'll need for French words such as {\it r\^ ole\/} and
{\it \'el\`eve}, for German words such as {\it Schu\ss}, and for words in several other languages
as well.  You'll find a complete list of \TeX's accents and letters of European languages in the
\TeX book. 

You can also get Greek letters such as ``$\alpha$'' and ``$\Omega$'' for use in math, card suits
such as ``$\spadesuit$'' and ``$\diamondsuit$'', music symbols such as ``$\sharp$'' and
``$\flat$'', and many other special symbols.  \TeX\ will only accept these sorts of special symbols
in its ``math mode'', so you'll to enclose them within `{\tt \$}' characters.  Here are some
more strange-looking symbols, $\Upsilon$, $\varpi$, $\partial$, $\wr$,
and $\hookleftarrow$.

Many files are involved in the \TeX\ system. The {\tt .mf} file contains the source codes for
the metafont program, which generate {\tt .tfm} and {\tt .pk} font files.  Most part of \TeX's
definitions are in the file {\tt plain.tex}.  
